\subsection{Half-Life} 
A common concept in exponential decay (that is, decrease), especially common in science, is the \emph{half-life}
\begin{definition}[Half-Life]
The \term{half-life} of a substance is the amount of time required for half of the substance to remain.
\end{definition}
\begin{example}
A certain radioactive isotope has a half-life of 183 days. A 50 gram sample is obtained. Write an exponential function describing the mass $m$ of the isotope that will remain after $t$ days.
\begin{itemize}
\item The starting value is \makeblank[1in] at \makeblank[1in]
\item The multiplier is \makeblank[1in]
\item $\Delta t$, the time it takes to multiply once, is \makeblanksmall
\end{itemize}
\[
m(t)=\makeblank
\]
\end{example}
We can also answer various questions.
\begin{example}
A certain radioactive isotope has a half-life of 183 days. A 50 gram sample is obtained. How much will remain after one year (365 days)? Round your answer to two decimal places.
\begin{lpic}{}{3}
\end{lpic}
\end{example}
We might also be asked to compute a half-life, after an exponential model is described in some other way.
\begin{example}
A drug is metabolized by the body in such a way that the amount of drug in the patient's system dereases by 22\% every hour. What is the half-life of the drug? Round to two decimal places.
\begin{lpic}{}{6}
\end{lpic}
\end{example}